%------------------------------------------------------------------------------%
\begin{question}
  Encontre o ponto da superfície \(z=xy+1\) que está mais próximo da origem
\end{question}

%------------------------------------------------------------------------------%
\begin{resolution}{Formulação do Problema}
  Queremos minimizar a distância até a origem
  \[
    d(x, y, z) = \sqrt{x^2 + y^2 + z^2}
  \]
  \pause
  Como a função raiz quadrada é crescente, podemos minimizar a função
  \[
    f(x, y, z) = x^2 + y^2 + z^2
  \]
  \pause
  sujeito a restrição
  \[
    g(x, y, z) = z - xy = 1
  \]
\end{resolution}

%------------------------------------------------------------------------------%
\begin{resolution}{Gradientes}
\begin{columns}
\column{0.5\linewidth}
  \onslide<+->{Gradiente de \(f(x, y, z) = x^2 + y^2 + z^2\)}
  \begin{align*}
    \onslide<+->{f_x(x, y, z) & = 2x \\}
    \onslide<+->{f_y(x, y, z) & = 2y \\}
    \onslide<+->{f_z(x, y, z) & = 2z}
  \end{align*}

  \onslide<+->{
  \[
    \nabla f = \vetortri{2x}{2y}{2z}
  \]}
\column{0.5\linewidth}
  \onslide<+->{Gradiente de \(g(x, y, z)=z-xy\)}
  \begin{align*}
    \onslide<+->{g_x(x, y, z) & = -y \\}
    \onslide<+->{g_y(x, y, z) & = -x \\}
    \onslide<+->{g_z(x, y, z) & = 1}
  \end{align*}
  \onslide<+->{
  \[
    \nabla g = \vetortri{-y}{-x}{1}
  \]}
\end{columns}
\end{resolution}

%------------------------------------------------------------------------------%
\begin{resolution}{Sistema}
  Aplicando os Multiplicadores de Lagrange, precisamos resolver o sistema
  \[
  \begin{cases}
    2x = -\lambda y \\
    2y = -\lambda x \\
    2z = \lambda  \\
    z-xy = 1
  \end{cases}
  \]
\end{resolution}

%------------------------------------------------------------------------------%
\begin{resolution}{Solução}
  Isolamos $x$ na primeira equação e substituímos na segunda
  \[
    \pause
    2x = -\lambda y
    \pause
    \hspace{54mm}
    x = \frac{-\lambda y}{2}
  \]
  \pause
  \[
    y
    \pause = \frac{-\lambda}{2} x
    \pause = \frac{-\lambda}{2} \; \frac{-\lambda y}{2}
    \pause = \frac{\lambda^2}{4} y
    \hspace{20mm}
    \pause 4y = \lambda^2 y
  \]

  \pause
  \vspace{\baselineskip}
  Temos dois casos \(y=0\) ou \(y\neq 0\)
\end{resolution}

%------------------------------------------------------------------------------%
\begin{resolution}{Caso 1}
  Caso \(y=0\)
  \[
    \pause
    x = \frac{-\lambda y}{2}
    \pause
    = 0
  \]

  \pause
  \(z-xy = 1\) se reduz a \(z=1\)

  \pause
  \vspace{\baselineskip}
  Obtemos o ponto \((x_1, y_1, z_1) = (0, 0, 1)\)
\end{resolution}

%------------------------------------------------------------------------------%
\begin{resolution}{Caso 2}
  Caso \(y\neq 0\)
  \[
    \pause
    4y = \lambda^2 y
  \]
  \[
    \pause
    \lambda^2 = 4
  \]
  \[
    \pause
    \lambda = \pm 2
  \]
\end{resolution}

%------------------------------------------------------------------------------%
\begin{resolution}{Caso 2}
  \begin{columns}[t]
  \column{0.5\linewidth}
    Se $y\neq0$ e $\lambda=2$
    \[
      \pause x
      \pause = \frac{-\lambda y}{2}
      \pause = \frac{-2y}{2}
      \pause = -y
    \]

    \[
      \pause 2z = \lambda
      \pause = 2
    \]
    \[
      \pause z = 1
    \]
  \column{0.5\linewidth}
    \pause
    Substituindo na quarta equação
    \pause
    \begin{align*}
      \onslide<+->{z-xy      & = 1 \\}
      \onslide<+->{1 - (-y)y & = 1 \\}
      \onslide<+->{y^2       & = 0 \\}
      \onslide<+->{y         & = 0   }
    \end{align*}
    \pause
    Contradição
  \end{columns}
\end{resolution}

%------------------------------------------------------------------------------%
\begin{resolution}{Caso 2}
  \begin{columns}[t]
  \column{0.5\linewidth}
    Se $y\neq 0$ e $\lambda=-2$
    \[
      \pause x
      \pause = \frac{-\lambda y}{2}
      \pause = \frac{-(-2)y}{2}
      \pause = y
    \]

    \[
      \pause 2z = \lambda = -2
    \]
    \[
      \pause z = -1
    \]
  \column{0.5\linewidth}
    \pause
    Substituindo na quarta equação
    \pause
    \begin{align*}
      \onslide<+->{z-xy     & = 1  \\}
      \onslide<+->{-1 - y^2 & = 1  \\}
      \onslide<+->{-y^2     & = 2  \\}
      \onslide<+->{y^2      & = -2}
    \end{align*}
  \end{columns}

  \vspace{\baselineskip}
  \onslide<+->{Não existe solução}
\end{resolution}

%------------------------------------------------------------------------------%
\begin{resolution}{Solução}
  O ponto mais próximo da origem é \((0, 0, 1)\)
\end{resolution}

%------------------------------------------------------------------------------%
